\documentclass[11pt,a4paper]{article}
\usepackage[utf8]{inputenc}
\usepackage[T1]{fontenc}
\usepackage{lmodern}
\usepackage[french]{babel}
\usepackage{amsmath,amssymb,mathtools}
\usepackage{icomma}
\usepackage{xcolor}
\usepackage[most]{tcolorbox}
\usepackage{enumitem}
\usepackage[a4paper,margin=2.2cm,headheight=26pt,headsep=10pt]{geometry}
\usepackage{fancyhdr}
\usepackage[hidelinks]{hyperref}
\definecolor{primary}{RGB}{222,49,99}
\definecolor{primarydark}{RGB}{150,22,55}
\tcbset{boxbase/.style={enhanced,breakable,boxrule=0.9pt,arc=2.5pt,
  left=3mm,right=3mm,top=2.3mm,bottom=2.3mm,fonttitle=\bfseries,coltitle=white}}
\newtcolorbox[auto counter]{exo}[1][]{boxbase,colback=primary!4,colframe=primary,
  title={\textsc{Exercice}~\thetcbcounter\ifx\relax#1\relax\relax\else\textmd{ --- \itshape #1}\fi}}
\newcommand{\R}{\mathbb{R}}
\newcommand{\ds}{\displaystyle}
\pagestyle{fancy}\fancyhf{}
\fancyhead[L]{\small\textcolor{primarydark}{\textbf{Thème 2} — Logarithmes}}
\fancyhead[R]{\small\textcolor{primarydark}{\textsc{Difficile}}}
\fancyfoot[C]{\small\thepage}
\renewcommand{\headrulewidth}{0.8pt}
\begin{document}

\begin{center}
{\LARGE\bfseries\color{primarydark} Niveau Difficile — Exercices}\\[2pt]
{\color{primary}\rule{0.45\textwidth}{1.2pt}}\\[3pt]
{\small\itshape Niveau concours. Les sous-questions construisent la difficulté pas à pas.}
\end{center}
\vspace{1mm}

\begin{exo}[manipuler un logarithme d'expression composée]
Soit $a,b,c$ trois réels strictement positifs et $\ds R=\frac{a^{2}b}{c}$.
\begin{enumerate}[label=\textbf{\alph*)},leftmargin=2.2em,topsep=1pt,itemsep=2pt]
  \item Exprime $\ln R$ en fonction de $\ln a$, $\ln b$ et $\ln c$.
  \item Montre que $\ln R-\ln b=2\ln a-\ln c$.
  \item Simplifie $\ln(Rc)-2\ln a$.
\end{enumerate}
\end{exo}

\begin{exo}[un carré parfait caché dans un logarithme]
On considère la fonction $f$ définie sur $\R$ par $f(x)=\ln\bigl(2e^{x}+1+e^{2x}\bigr)$.
\begin{enumerate}[label=\textbf{\alph*)},leftmargin=2.2em,topsep=1pt,itemsep=2pt]
  \item Montre que l'argument du logarithme est un carré parfait, et factorise-le.
  \item Déduis-en que $f(x)=2\ln\bigl(1+e^{x}\bigr)$.
  \item En factorisant $e^{x}$ à l'intérieur de $1+e^{x}$, montre que $f(x)=2\bigl[x+\ln\bigl(1+e^{-x}\bigr)\bigr]$.
\end{enumerate}
\end{exo}

\begin{exo}[mise en situation : croissance exponentielle]
Une culture de cellules croît de façon exponentielle : le nombre de cellules \emph{double tous les
$20$~jours}. On note $N_0$ le nombre initial et $N(t)$ le nombre après $t$ jours.
\begin{enumerate}[label=\textbf{\alph*)},leftmargin=2.2em,topsep=1pt,itemsep=2pt]
  \item Justifie que l'on peut écrire $\ds N(t)=N_0\times 2^{\,t/20}$.
  \item Après combien de jours le nombre de cellules est-il \emph{multiplié par $8$}~? (Ramène à une même base, sans logarithme.)
  \item On veut maintenant multiplier le nombre initial par $10$. Exprime le temps $t$ nécessaire à l'aide d'un logarithme. Donne aussi le facteur par lequel la population est multipliée \emph{en un seul jour}, sous la forme d'une puissance de $2$.
\end{enumerate}
\end{exo}

\end{document}
